\paragraph{Leitura.} As ablações mostram que a decisão depende de duas regras: a ativação de TS2 (colocalização multimétrica) e a do modificador M1 (granularidade do bucket, que materializa TS2-U). Sem qualquer uma delas, o HBench perde o melhor deployment e o regret medido sobe a 31,6\%, pior que a escolha aleatória (6,4\% esperado). As regras de TS3, TS6, TS7, TS8, M2, M4 e M5 não ativam nesta carga, então não se pode afirmar nada sobre a contribuição delas aqui; isso reforça a limitação de diversidade do workload. Sem Pareto, sem diversidade ou com $k=1$ o melhor também é encontrado, porque o escore ponderado já põe TS2-U em primeiro.
